% يجب الترجمة بمحرك LuaLaTeX  (Menu -> Compiler -> LuaLaTeX)
\documentclass[12pt, a4paper]{article}
\usepackage[a4paper, top=1.2cm, bottom=1.5cm, left=1.5cm, right=1.5cm]{geometry}
\usepackage{amsmath, amssymb}
\usepackage{fontspec}
\usepackage[bidi=basic]{babel}
\babelprovide[main, import, onchar=ids fonts]{arabic}
\babelprovide[import, onchar=ids fonts]{english}
\babelfont{rm}{Noto Serif}
\babelfont[arabic]{rm}{Amiri} % الخط الأميري
\usepackage{enumitem}
\usepackage{tabularx}
\usepackage{array}
\usepackage{graphicx}
\usepackage{tikz}
\usetikzlibrary{calc, arrows.meta}
\usepackage{eso-pic}

% الصفحة عربية (من اليمين لليسار)، فنضع كل رسمة داخل صندوق إنجليزي
% حتى لا تنعكس أماكن الأرقام والرموز داخل الرسومات
\newcommand{\LTRfig}[1]{\mbox{\foreignlanguage{english}{#1}}}

% صف خيارات من أربعة أعمدة
\newcommand{\choices}[4]{%
\vspace{0.15cm}\par\noindent
\begin{tabularx}{\linewidth}{XXXX}
(a) #1 & (b) #2 & (c) #3 & (d) #4
\end{tabularx}}

% مستوى إحداثي فارغ للرسم عليه
\newcommand{\blankplane}{%
\LTRfig{%
\begin{tikzpicture}[scale=0.38]
  \draw[gray!35, very thin] (-7,-5) grid (7,5);
  \draw[-{Stealth}, thick] (-7.4,0) -- (7.6,0) node[right] {$x$};
  \draw[-{Stealth}, thick] (0,-5.4) -- (0,5.6) node[above] {$y$};
  \foreach \i in {-6,-4,-2,2,4,6} \node[below, font=\tiny] at (\i,0) {$\i$};
  \foreach \i in {-4,-2,2,4} \node[left, font=\tiny] at (0,\i) {$\i$};
\end{tikzpicture}}}

\begin{document}

% إطار الصفحات
\AddToShipoutPictureBG{%
  \LTRfig{%
  \begin{tikzpicture}[remember picture, overlay]
    \draw[line width=1.2pt] ([shift={(0.8cm,-0.8cm)}]current page.north west) rectangle ([shift={(-0.8cm,0.8cm)}]current page.south east);
  \end{tikzpicture}}%
}

% الترويسة
\noindent
\begin{minipage}[t]{0.38\textwidth}
\vspace{0pt}
\textbf{المملكة الأردنية الهاشمية} \\
\textbf{وزارة التربية والتعليم} \\
\textbf{مدرسة عنبة الثانوية الشاملة للبنات} \\
\textbf{المبحث: الرياضيات}
\end{minipage}%
\begin{minipage}[t]{0.24\textwidth}
\vspace{0pt}
\centering
\IfFileExists{logo.png}{\includegraphics[width=2.3cm]{logo.png}}{}
\end{minipage}%
\begin{minipage}[t]{0.38\textwidth}
\vspace{0pt}
\textbf{الصف: التاسع} \\
\textbf{الشعبة: ( \hspace{1.5cm} )} \\
\textbf{اليوم والتاريخ: \hspace{0.5cm} / \hspace{0.5cm} / 202 م} \\
\textbf{الزمن: \hspace{2cm}}
\end{minipage}

\vspace{0.3cm}
\hrule height 1.2pt
\vspace{0.08cm}
\hrule height 0.4pt
\vspace{0.3cm}

\begin{center}
\textbf{\Large اختبار الرياضيات}
\end{center}

\noindent \textbf{اسم الطالبة:} \rule{10cm}{0.5pt}

\vspace{0.4cm}
\noindent\textbf{\large السؤال الأول: اختر الإجابة الصحيحة في كل مما يأتي: \hfill (20 علامة)}
\vspace{0.1cm}

\begin{enumerate}[label=\textbf{\arabic*)}, itemsep=0.25cm, leftmargin=*]

\item حل المتباينة $-9x + 17 \geq -64$ هو:
\choices{$\{x \mid x \leq -9\}$}{$\{x \mid x \geq 9\}$}{$\{x \mid x \leq 9\}$}{$\{x \mid x \geq -9\}$}

\item الفترة التي تعبّر عن التمثيل البياني الآتي هي: \hfill
\raisebox{-0.25cm}{\LTRfig{%
\begin{tikzpicture}[scale=0.8]
  \draw[{Stealth}-{Stealth}, thick] (1.3,0) -- (6.8,0);
  \foreach \i in {2,...,6} {\draw (\i,0.1) -- (\i,-0.1); \node[below, font=\small] at (\i,-0.05) {$\i$};}
  \draw[-{Stealth}, very thick, blue!70!black] (4,0.25) -- (6.8,0.25);
  \fill[blue!70!black] (4,0.25) circle (3.5pt);
\end{tikzpicture}}}
\choices{$(4, \infty)$}{$[4, \infty)$}{$(-\infty, 4)$}{$(-\infty, 4]$}

\item الفترة التي تمثل حل المتباينة $1 - x \geq 6$ هي:
\choices{$(-\infty, -5]$}{$(5, \infty)$}{$[-5, \infty)$}{$(-\infty, -5)$}

\item أي الأزواج الآتية يعتبر أحد حلول المتباينة $4y + x < 5$ ؟
\choices{$(2, 1)$}{$(1, 1)$}{$(1, -1)$}{$(-2, 2)$}

\item إذا كانت $B$ مجموعة الأعداد الكلية الأقل من $1$، فإن نوع المجموعة $B$ هو:
\choices{خالية}{مفردة}{منتهية}{غير منتهية}

\item المتباينة التي تمثل الجملة: «عدد على الأكثر $-4$ أو على الأقل $2$» هي:
\vspace{0.15cm}\par\noindent
\begin{tabularx}{\linewidth}{XX}
(a) $x \leq 2 \ \text{or}\ x \geq -4$ & (b) $x < 2 \ \text{or}\ x > -4$ \\[0.1cm]
(c) $x \geq 2 \ \text{or}\ x \leq -3$ & (d) $x > 2 \ \text{or}\ x < -3$
\end{tabularx}

\item مجال العلاقة الآتية $\{(-3,-2),\ (4,1),\ (2,3),\ (1,4),\ (0,2)\}$ هو:
\choices{$\{-2, 3, 4\}$}{$\{-2, 1, 2, 4\}$}{$\{-3, 0, 1, 2\}$}{$\{-3, 0, 1, 2, 4\}$}

\item مجموعة حل المتباينة $-7 < x + 2 < 4$ هي:
\choices{$(-5, 6)$}{$(-9, 6)$}{$(-5, 2)$}{$(-9, 2)$}

\item إحداثيا نقطة رأس القطع المكافئ للاقتران التربيعي $y = x^2 + 2x + 3$ هما:
\choices{$(0, 3)$}{$(2, 11)$}{$(-1, 2)$}{$(1, 6)$}

\item مدى الاقتران التربيعي $f(x) = 12x - 3x^2 + 3$ هو:
\choices{$\{y \mid y \leq 15\}$}{$\{y \mid y \leq 3\}$}{$\{y \mid y \geq 15\}$}{$\{y \mid y \geq 3\}$}

\end{enumerate}

\newpage
% ===================== الصفحة الثانية: السؤالان الثاني والثالث =====================
% (الأرقام بجانب \vspace هي مساحة الحل، ويمكن تكبيرها أو تصغيرها)
\noindent\textbf{\large السؤال الثاني:}
\vspace{0.3cm}

\noindent\textbf{أ) أكتب ما يأتي بطريقة سرد العناصر، وطريقة الصفة المميزة: \hfill (4 علامات)}
\begin{enumerate}[label=\textbf{\arabic*)}, leftmargin=*]
\item مجموعة الأعداد الزوجية التي تقل عن أو تساوي $16$.
\vspace{2.2cm} % سطران: سرد العناصر + الصفة المميزة
\item مجموعة حل المعادلة $5x - 30 = 0$
\vspace{2.2cm} % حل المعادلة + الطريقتان
\end{enumerate}

\noindent\textbf{ب) أكتب مجموعة حل المتباينة باستعمال الصفة المميزة، ثم أمثلها على خط الأعداد: \hfill (4 علامات)}
\begin{enumerate}[label=\textbf{\arabic*)}, leftmargin=*]
\item $3x - 8 \leq 10$
\end{enumerate}
\vspace{2.2cm} % ثلاث خطوات حل + مجموعة الحل
\begin{center}
\LTRfig{%
\begin{tikzpicture}[scale=0.85]
  \draw[{Stealth}-{Stealth}, thick] (-0.6,0) -- (10.6,0);
  \foreach \i in {0,...,10} \draw (\i,0.12) -- (\i,-0.12);
\end{tikzpicture}}
\end{center}

\vspace{0.3cm}
\hrule height 0.4pt
\vspace{0.5cm}

\noindent\textbf{\large السؤال الثالث: أجد مجموعة حل كل من المتباينات الآتية، وأكتب مجموعة الحل باستعمال رمز الفترة: \hfill \mbox{(6 علامات)}}
\begin{enumerate}[label=\textbf{\arabic*)}, leftmargin=*]
\item $-2 \leq x - 7 \leq 1$
\vspace{3.8cm} % متباينة مركبة: سطران أو ثلاثة + الفترة
\item $2x + 3 < 7 \quad \text{or} \quad x + 9 \leq 2$
\vspace{3.5cm} % متباينتان + الاتحاد + الفترة
\end{enumerate}

\newpage
% ===================== الصفحة الثالثة: السؤال الرابع كاملاً =====================
\noindent\textbf{\large السؤال الرابع:}
\vspace{0.3cm}

\noindent\textbf{أ) أمثل كلاً من المتباينات الآتية في المستوى الإحداثي: \hfill (7 علامات)}
\vspace{0.2cm}

\noindent\begin{tabularx}{\linewidth}{XX}
\textbf{1)} \ $4y - 8 \geq 0$ & \textbf{2)} \ $x < 6$ \\[0.2cm]
\blankplane & \blankplane
\end{tabularx}

\vspace{0.4cm}
\hrule height 0.4pt
\vspace{0.4cm}

\noindent\textbf{ب) أجد مجال ومدى كل مما يأتي، وأحدد إذا كان يمثل اقتراناً متصلاً أو منفصلاً: \hfill (5 علامات)}
\begin{enumerate}[label=\textbf{\arabic*)}, leftmargin=*, itemsep=0.2cm]
\item $\{(4,2),\ (1,1),\ (0,0),\ (1,-1),\ (4,5)\}$
\vspace{2.6cm} % المجال + المدى + نوعه
\item % الرسم على اليسار، ومساحة الإجابة على اليمين بجانبه
\hfill
\begin{minipage}[t]{0.5\linewidth}
\vspace{0pt}
\raggedright
\LTRfig{%
\begin{tikzpicture}[scale=0.5]
  \draw[gray!30, very thin] (-3,-4) grid (7,4);
  \draw[-{Stealth}, thick] (-3.4,0) -- (7.5,0) node[right] {$x$};
  \draw[-{Stealth}, thick] (0,-4.4) -- (0,4.5) node[above] {$y$};
  \foreach \i in {-2,-1,1,2,3,4,5,6} \node[below, font=\scriptsize, fill=white, inner sep=1pt] at (\i,-0.1) {$\i$};
  \foreach \i in {-3,-2,-1,1,2,3} \node[right, font=\scriptsize, fill=white, inner sep=1pt] at (0.05,\i) {$\i$};
  % y = 5/16 (x-2)^2 - 3 , يبدأ من النقطة (-2, 2)
  \draw[very thick, blue!70!black, -{Stealth}, domain=-2:6.55, samples=80]
       plot (\x, {5/16*(\x-2)^2 - 3});
  \fill[blue!70!black] (-2,2) circle (4pt);
  \draw[dashed] (-2,2) -- (0,2);
  \draw[dashed] (2,-3) -- (0,-3);
\end{tikzpicture}}
\end{minipage}
\end{enumerate}

\vspace{0.4cm}
\hrule height 0.4pt
\vspace{0.4cm}

\noindent\textbf{ج) إذا كان $g(x) = 16x - 2$ ، فأجد: \hfill (6 علامات)}
\vspace{0.2cm}

\noindent\begin{tabularx}{\linewidth}{XXX}
\textbf{1)} \ $g(1)$ & \textbf{2)} \ $2\,g(2) - 1$ & \textbf{3)} \ قيمة $x$ التي تجعل $g(x) = 30$
\end{tabularx}
\vspace{2.8cm} % ثلاث عمليات قصيرة متجاورة

\newpage
% ===================== الصفحة الرابعة: السؤال الخامس =====================
\noindent\textbf{\large السؤال الخامس: يبيّن التمثيل البياني المجاور رحلة علي على دراجته من منزله إلى المسجد، حيث أمضى بعض الوقت فيه، ثم عاد بدراجته إلى المنزل. \hfill \mbox{(8 علامات)}}
\vspace{0.3cm}

\noindent
\begin{minipage}[t]{0.47\textwidth}
\vspace{0pt}
\begin{enumerate}[label=\textbf{\alph*)}, itemsep=1.5cm, leftmargin=*] % سطر إجابة لكل من أ، ب، ج
\item[\textbf{أ)}] ما البعد بين منزل علي والمسجد؟
\item[\textbf{ب)}] في أي ساعة غادر علي منزله؟
\item[\textbf{ج)}] كم أمضى علي من الوقت في المسجد؟
\item[\textbf{د)}] أجد السرعة المتوسطة لعلي في المدة الزمنية من \mbox{$10{:}00$} إلى \mbox{$10{:}30$}، ثم أبيّن ماذا تمثل.
\end{enumerate}
\vspace{3cm} % د) قانون السرعة + التعويض + التفسير
\end{minipage}%
\hfill
\begin{minipage}[t]{0.5\textwidth}
\vspace{0pt}
\centering
\LTRfig{%
\begin{tikzpicture}[xscale=1.05, yscale=0.65]
  \draw[gray!30, very thin] (0,0) grid[xstep=0.5, ystep=1] (6,6);
  \draw[-{Stealth}, thick] (0,0) -- (6.5,0) node[right] {\foreignlanguage{arabic}{الزمن}};
  \draw[-{Stealth}, thick] (0,0) -- (0,6.8) node[above, align=center] {\foreignlanguage{arabic}{المسافة}\\[-2pt] $(\mathrm{km})$};
  \foreach \t/\lab in {0/8:00, 1/8:30, 2/9:00, 3/9:30, 4/10:00, 5/10:30, 6/11:00}
     \node[below, font=\scriptsize, rotate=45, anchor=north east] at (\t,0) {\lab};
  \foreach \d in {1,...,6} \node[left, font=\scriptsize] at (0,\d) {$\d$};
  % رحلة علي
  \draw[very thick, blue!70!black] (0,0) -- (2,6) -- (4,6) -- (5,0);
  \fill[blue!70!black] (0,0) circle (2pt) (2,6) circle (2pt) (4,6) circle (2pt) (5,0) circle (2pt);
\end{tikzpicture}}
\end{minipage}

\vfill
\begin{center}
\textbf{\large مع تمنياتي لكم بالتوفيق دائماً} \\
\vspace{0.3cm}
\textbf{معلمتكم: ولاء أبو شرية}
\end{center}

\end{document}
